\section{Spatial Cloaking for Movement Ecology}


\subsection{Problem Formulation}

Given a trajectory $\tau = \{(\mathbf{x}_1, t_1), (\mathbf{x}_2, t_2), \ldots, (\mathbf{x}_n, t_n)\}$ where $\mathbf{x}_i \in \mathbb{R}^2$ are spatial coordinates and $t_i$ are timestamps, we seek a transformation $\mathcal{T}: \tau \to \tilde{\tau}$ that satisfies spatial $k$-anonymity while minimizing movement metric distortion.

\subsection{Habitat-Aware Cloaking}

Na\"ive spatial cloaking (e.g., adding uniform noise) destroys ecological validity by placing cloaked locations in impossible habitats. We introduce \emph{habitat-aware cloaking} that constrains perturbations to ecologically plausible regions.

\begin{definition}[Habitat Suitability Mask]
For species $s$, the habitat suitability mask $H_s: \mathbb{R}^2 \to [0, 1]$ assigns each location a probability of suitability based on land cover, elevation, proximity to water, and known range boundaries.
\end{definition}

The cloaking algorithm perturbs each location $\mathbf{x}_i$ to $\tilde{\mathbf{x}}_i$ drawn from:
\begin{equation}
\tilde{\mathbf{x}}_i \sim \frac{H_s(\mathbf{x}) \cdot \mathcal{N}(\mathbf{x}_i, \sigma^2 \mathbf{I})}{\int H_s(\mathbf{x}') \cdot \mathcal{N}(\mathbf{x}'; \mathbf{x}_i, \sigma^2 \mathbf{I}) \, d\mathbf{x}'}
\end{equation}
where $\sigma$ is calibrated to achieve the required cloaking radius $R_{\text{cloak}}$ for the species' privacy tier.

\subsection{Trajectory-Consistent Cloaking}

Independent perturbation of each location introduces artificial discontinuities in movement trajectories. We enforce trajectory consistency through a correlated noise process:

\begin{equation}
\boldsymbol{\eta}_i = \rho \cdot \boldsymbol{\eta}_{i-1} + \sqrt{1 - \rho^2} \cdot \boldsymbol{\xi}_i, \quad \boldsymbol{\xi}_i \sim \mathcal{N}(\mathbf{0}, \sigma^2 \mathbf{I})
\end{equation}

where $\rho \in [0, 1)$ is the autocorrelation parameter that controls the smoothness of the perturbation trajectory, and $\tilde{\mathbf{x}}_i = \mathbf{x}_i + \boldsymbol{\eta}_i$.

\begin{theorem}[Movement Metric Preservation]
\label{thm:preservation}
Under trajectory-consistent habitat-aware cloaking with autocorrelation $\rho$ and noise scale $\sigma$, the step length distribution $\ell$ and turning angle distribution $\theta$ of the cloaked trajectory $\tilde{\tau}$ satisfy:
\begin{align}
\mathbb{E}[|\bar{\ell}(\tilde{\tau}) - \bar{\ell}(\tau)|] &\leq \sigma\sqrt{2(1-\rho)} \\
\mathbb{E}[|\text{Var}(\theta(\tilde{\tau})) - \text{Var}(\theta(\tau))|] &\leq \frac{4\sigma^2(1-\rho)}{(\bar{\ell}(\tau))^2}
\end{align}
where $\bar{\ell}(\tau)$ is the mean step length of trajectory $\tau$.
\end{theorem}

\begin{proof}
The cloaked step vector is $\tilde{\mathbf{d}}_i = \tilde{\mathbf{x}}_{i+1} - \tilde{\mathbf{x}}_i = \mathbf{d}_i + (\boldsymbol{\eta}_{i+1} - \boldsymbol{\eta}_i)$. The perturbation difference $\boldsymbol{\eta}_{i+1} - \boldsymbol{\eta}_i = (\rho - 1)\boldsymbol{\eta}_i + \sqrt{1-\rho^2}\boldsymbol{\xi}_{i+1}$, which has variance $\sigma^2(2 - 2\rho)$ in each dimension. The expected absolute error in step length follows from the triangle inequality applied to the noise magnitude $\|\boldsymbol{\eta}_{i+1} - \boldsymbol{\eta}_i\|$, whose expected value is bounded by $\sigma\sqrt{2(1-\rho)} \cdot \sqrt{\pi/2}$. The turning angle bound follows from the small-angle approximation for the angular perturbation $\arctan(\sigma\sqrt{2(1-\rho)}/\bar{\ell})$ applied to both components of the turn.
\end{proof}

\subsection{Home Range Estimation Under Cloaking}

Home range estimation is a fundamental ecological metric. We show that kernel density estimation (KDE) home ranges are robust to our cloaking algorithm.

\begin{proposition}[Home Range Bound]
Let $\hat{A}_\alpha$ and $\tilde{A}_\alpha$ be the $\alpha$-level home range areas estimated from original and cloaked trajectories using KDE with bandwidth $h$. Then:
\begin{equation}
|\tilde{A}_\alpha - \hat{A}_\alpha| \leq 2\pi \sigma^2 + 4\sigma h \sqrt{2\ln(1/\alpha)} + O(\sigma^3)
\end{equation}
\end{proposition}

For typical KDE bandwidths ($h \geq 500$m) and cloaking scales ($\sigma = 1$--5 km), this implies home range estimation errors below 8\% for 95\% home ranges.

